% !TEX root = ../main.tex
\chapter{Literature Review and Theoretical Framework}
\label{ch:literature}

Chapter~\ref{ch:introduction} stated the research problem---that transfer-graph reputation scores conflate payment flow with trust---and derived objectives O1--O4 and research questions RQ1--RQ4 from it. This chapter reviews the conceptual and technical foundations required to interpret EndorseRank and Adaptive Weighted PageRank (AWP) and to judge whether those objectives are attainable with public data. It begins with blockchain and DeFi primitives for readers who may not specialize in distributed ledgers, then develops graph ranking theory, reputation-system lineage, ERC-20 allowance semantics, and rank-correlation statistics. A theoretical framework---hyperlink--citation propagation, domain alignment, and computational complexity---motivates the primary EndorseRank versus AWP comparison. Later sections supply a terminology guide, worked example, risk taxonomy, AWP in depth, a tabular summary of the reviewed literature (methodology, strengths, limitations), and an explicit statement of the research gap that Chapters~\ref{ch:methodology} and~\ref{ch:results} address.

\dissertationSubheading[sec:blockchain-defi-foundations]{Blockchain and DeFi foundations}

A blockchain is a replicated append-only ledger in which batches of transactions (blocks) are linked by cryptographic hashes \cend{5}. Participants agree on a canonical history through a consensus mechanism; once confirmed, entries are costly to rewrite. This property enables \emph{trust-minimized} settlement: counterparties need not share a legal identity if they can verify the ledger state.

Bitcoin demonstrated peer-to-peer electronic cash \cend{5}. Ethereum generalized the model to a world computer: accounts hold balances and code, and transactions invoke functions that update state \cend{7}; the execution semantics are specified by the Ethereum Virtual Machine (EVM) \cend{36}. Smart contracts are programs stored on-chain that execute deterministically when called \cend{6}. Because execution is public and reproducible, financial logic---collateral checks, liquidations, fee accrual---can be encoded without a central operator \cend{1}.

Decentralized Finance (DeFi) refers to the ecosystem of such contracts offering exchange, lending, derivatives, and asset management; Werner et al.\cend{8} give a systematization of the protocol classes, and Sch\"{a}r\cend{1} describes the layered architecture. Composability---the ability of one protocol to call another---creates powerful products \cend{37} but also systemic risk pathways: flash-loan attacks \cend{38} and price-oracle manipulation \cend{39} are the canonical examples. For reputation research, the relevant implication is that \emph{wallet-level histories} on public chains are rich, machine-readable, and incomplete: they reveal actions, not legal identities \cend{2}.

\dissertationSubsubheading[sub:pseudonymity-credit]{Pseudonymity and the credit problem}

Traditional credit scoring links legal persons to repayment histories. On public blockchains, the unit of analysis is an address (wallet). An address may represent an individual, a fund, a bot, or a smart contract. Multiple addresses may be controlled by one entity \cend{11}, and clustering heuristics recover such control only partially \cend{40}. Consequently, DeFi lending markets typically require over-collateralization: borrowers lock assets worth more than the loan, and liquidators seize collateral when health factors fall \cend{41}. Over-collateralization is robust but capital-inefficient. On-chain reputation scores aim to summarize behavioral reliability so that protocols---or researchers studying them---can reason about risk without off-chain KYC; proposals range from prospect-theory credit models \cend{13} to machine-learning scoring on lending data \cend{14} and risk propagation over account networks \cend{12}.

\dissertationSubheading[sec:ethereum-l2]{Ethereum, Layer-2 rollups, and Arbitrum One}

Ethereum mainnet provides security and liquidity but faces throughput and fee constraints under demand \cend{42}. Layer-2 (L2) protocols execute transactions off the base layer (or in a compressed form) and periodically post commitments to Ethereum, aiming to inherit security while reducing cost \cend{43}; Thibault et al.\cend{44} survey the rollup design space. Optimistic rollups such as Arbitrum assume correctness unless challenged during a dispute window \cend{45}.

Arbitrum One is an EVM-compatible optimistic rollup widely used for DeFi activity, including the GMX perpetual-swap venue \cend{46} whose contracts appear among the most-approved spenders in the data; perpetual futures themselves are a distinct derivative class whose funding-rate pricing is analysed by Ackerer et al.\cend{47}. For this research, Arbitrum One is the empirical setting because (i) ERC-20 \texttt{Approval} and \texttt{Transfer} logs are available in public BigQuery datasets and (ii) fee levels permit dense activity within a six-month window. Results should not be assumed to transfer automatically to other L2s or to Ethereum mainnet without re-validation, but the methodological template---allowance graph versus transfer graph, proxy families, rank correlation---is chain-agnostic wherever equivalent logs exist.

\dissertationSubheading[sec:erc20-allowance-lifecycle]{ERC-20 tokens and the allowance lifecycle}

The ERC-20 standard defines a common interface for fungible tokens: \texttt{transfer}, \texttt{approve}, \texttt{transferFrom}, \texttt{allowance}, and related events \cend{18}. When an owner calls \texttt{approve(spender, amount)}, the token contract records that \texttt{spender} may transfer up to \texttt{amount} from the owner. The spender later calls \texttt{transferFrom} to move tokens. EIP-2612 \texttt{permit} allows allowance updates via off-chain signatures settled in a single transaction \cend{19}, improving user experience without changing the fundamental authorization semantics.

Figure~\ref{fig:allowance-sequence} illustrates the lifecycle. An allowance is an \emph{explicit authorization}, not a completed transfer. Owners may set allowances to routers, vaults, or trading contracts they trust to act within limits. Revocation (approving zero) and replacement (a new non-zero allowance) update the live authorization state. EndorseRank therefore uses the \emph{latest} allowance per owner--spender pair as edge strength: older approvals are superseded rather than decayed continuously.

\begin{figure}[htbp]
\centering
\begin{tikzpicture}[
  font=\small,
  actor/.style={draw, rounded corners, minimum width=2.4cm, minimum height=0.75cm, align=center},
  msg/.style={-{Stealth[length=2.5mm]}, thick},
  note/.style={font=\scriptsize\itshape, text=black!70, align=left}
]
% Actors
\node[actor] (owner) at (0,0) {Owner};
\node[actor] (token) at (5.5,0) {ERC-20 token};
\node[actor] (spender) at (11,0) {Spender};

% Lifelines (vertical dashed lines)
\draw[dashed, gray] (owner.south) -- ++(0,-5.2);
\draw[dashed, gray] (token.south) -- ++(0,-5.2);
\draw[dashed, gray] (spender.south) -- ++(0,-5.2);

% Message 1: approve (owner -> token), horizontal span 5.5cm
\draw[msg] (0,-1.2) -- node[above,font=\scriptsize] {\texttt{approve(spender,$a$)}} (5.5,-1.2);
\node[note] at (2.75,-1.7) {emit \texttt{Approval}; store allowance};

% Message 2: transferFrom (spender -> token), horizontal span 5.5cm
\draw[msg] (11,-2.8) -- node[above,font=\scriptsize] {\texttt{transferFrom(owner,to,$v\leq a$)}} (5.5,-2.8);
\node[note] at (8.25,-3.3) {move tokens; reduce allowance};

% Message 3: revoke (owner -> token), dashed
\draw[msg, dashed] (0,-4.4) -- node[above,font=\scriptsize] {\texttt{approve(spender,$0$)} (revoke)} (5.5,-4.4);
\end{tikzpicture}
\caption{ERC-20 allowance lifecycle: approve (or permit), optional \texttt{transferFrom}, and revoke. EndorseRank edges reflect the latest live allowance snapshot per owner--spender pair (historical approvals superseded, not time-decayed).}
\label{fig:allowance-sequence}
\end{figure}

From a measurement perspective, allowances are sparser than transfers: many transfers occur under a single approval, and many wallets never approve third-party spenders. This sparsity is both a computational advantage and a semantic claim---EndorseRank measures authorization networks, not payment networks.

\dissertationSubheading[sec:graph-reputation]{Graph-based on-chain reputation}

PageRank and variants rank nodes by propagated influence through directed edges \cend{15}. Two recent lines of work apply this idea to blockchain transaction graphs and must be kept distinct. Do and Do\cend{4} apply PageRank and HodgeRank to Ethereum transaction graphs as a social-credit measure; their contribution is the framing of transfer connectivity as credit-relevant. Adaptive Weighted PageRank (AWP), proposed by Do, Do, and Nguyen\cend{3}, is the transfer-graph baseline used in this thesis: it weights each transfer edge by value and by a logistic time-decay of edge age, and it validates the resulting scores against in-degree and in-value---count and sum of inbound transfers---as domain-intuition proxies for economic importance. The logistic decay formulation and the in-degree/in-value validation protocol reproduced in Chapter~\ref{ch:methodology} are therefore attributed to the AWP paper\cend{3}, not to Do and Do\cend{4}. Lin et al.\cend{12} show that transaction connectivity alone can miss risk-relevant structure unless risk signals propagate through the network, motivating refined graph semantics beyond vanilla PageRank; Wu et al.\cend{24} reach a similar conclusion for phishing detection, where network embeddings of transaction graphs outperform per-account features. Surveys of blockchain-based reputation systems \cend{25} and of smart-contract reputation designs \cend{26} document that most deployed systems still rely on explicit ratings or transfer volume as the trust signal; a survey of knowledge discovery in cryptocurrency transaction data \cend{48} likewise lists degree- and value-based centralities as the dominant wallet descriptors. Packin and Lev-Aretz\cend{2} emphasize auditability constraints under pseudonymity. EndorseRank stays within interpretable graph propagation. It is not a black-box predictive score.

Graph-based approaches are attractive because they (i) use only public data, (ii) produce continuous scores suitable for ranking, and (iii) admit complexity analysis in terms of edge counts. They are limited by Sybil attacks \cend{11}, wash trading, and the gap between graph centrality and economic default risk. PageRank scores are not treated as credit scores in the regulatory sense. The claim is that allowance-based and transfer-based rankings can be compared rigorously on shared proxies.

\dissertationSubheading[sec:pagerank-formalism]{PageRank formalism}

Let $G=(V,E)$ be a directed graph with $|V|=n$ nodes. Write $w_{ij}\geq 0$ for the weight of edge $i\to j$ (influence from $i$ toward $j$ in the endorsement or transfer sense used here). Let $W$ be the weighted adjacency matrix and define the out-strength $d_i^{\mathrm{out}}=\sum_j w_{ij}$. The row-stochastic transition matrix $P$ is
\begin{equation}
P_{ij} =
\begin{cases}
w_{ij}/d_i^{\mathrm{out}} & \text{if } d_i^{\mathrm{out}}>0,\\
1/n & \text{if } d_i^{\mathrm{out}}=0
\end{cases}
\label{eq:transition}
\end{equation}
(dangling nodes distribute mass uniformly). The Google matrix with damping $d\in(0,1)$ is
\begin{equation}
G_{\mathrm{PR}} = d\, P^{\top} + \frac{1-d}{n}\,\mathbf{1}\mathbf{1}^{\top},
\label{eq:google-matrix}
\end{equation}
and the PageRank vector $\boldsymbol{\pi}$ satisfies $\boldsymbol{\pi}=G_{\mathrm{PR}}\boldsymbol{\pi}$ with $\boldsymbol{\pi}\geq 0$ and $\sum_i \pi_i=1$ \cend{15,23}. Equivalently,
\begin{equation}
\pi_j = d\sum_{i:(i\to j)\in E}\pi_i \frac{w_{ij}}{d_i^{\mathrm{out}}} + \frac{1-d}{n}.
\label{eq:pagerank-node}
\end{equation}

\dissertationSubsubheading[sub:random-surfer]{Random-surfer interpretation}

Equation~\eqref{eq:pagerank-node} admits a probabilistic reading: a surfer follows outgoing edges with probability $d$ and teleports uniformly with probability $1-d$ (Figure~\ref{fig:random-surfer}). The stationary distribution ranks nodes by long-run visit frequency. Damping prevents sinks from trapping all mass and guarantees primitivity of $G_{\mathrm{PR}}$ under mild conditions, so the Perron vector is unique \cend{23}.

\begin{figure}[htbp]
\centering
\begin{tikzpicture}[
  font=\small,
  vertex/.style={draw, circle, minimum size=0.95cm},
  arrow/.style={-{Stealth[length=2.5mm]}, thick}
]
\node[vertex] (a) at (0,0) {$A$};
\node[vertex] (b) at (3.2,1.4) {$B$};
\node[vertex] (c) at (3.2,-1.4) {$C$};
\node[vertex] (d) at (6.4,0) {$D$};
\draw[arrow] (a) -- (b);
\draw[arrow] (a) -- (c);
\draw[arrow] (b) -- (d);
\draw[arrow] (c) -- (d);
\draw[arrow, dashed] (d) to[bend left=25] (a);
\node[align=left, font=\scriptsize] at (3.2,-2.8) {solid: follow edge (prob.\ $d$); dashed: teleport (prob.\ $1-d$)};
\end{tikzpicture}
\caption{Random-surfer intuition for PageRank. Solid arrows are graph transitions; dashed arrows represent teleportation that mixes the chain.}
\label{fig:random-surfer}
\end{figure}

\dissertationSubsubheading[sub:power-iteration]{Power iteration and complexity}

In practice $\boldsymbol{\pi}$ is computed by power iteration:
\begin{equation}
\boldsymbol{\pi}^{(t+1)} = G_{\mathrm{PR}}\boldsymbol{\pi}^{(t)},
\label{eq:power-iteration}
\end{equation}
until $\|\boldsymbol{\pi}^{(t+1)}-\boldsymbol{\pi}^{(t)}\|_1 < \varepsilon$ or a maximum iteration count is reached. Each iteration costs $O(|E|)$ arithmetic operations for sparse graphs \cend{49}. Therefore, methods that induce fewer edges---such as latest-allowance snapshots versus full transfer histories---reduce per-iteration cost and often memory. This complexity observation grounds hypothesis H3.

Default parameters in this research use damping $d=0.85$ following common practice \cend{15}, with convergence tolerance $\varepsilon=10^{-8}$ and a maximum of $300$ iterations as study-specific solver settings.\footnote{Tolerance and iteration caps are implementation choices for reproducible timing; they are not claimed as part of the classical PageRank formulation.} Robustness checks vary $d\in\{0.75,0.85,0.95\}$.

\dissertationSubheading[sec:reputation-systems]{Reputation systems lineage}

PageRank is one member of a family of spectral and propagation-based reputation mechanisms. Kleinberg's HITS algorithm distinguishes hubs and authorities in hyperlinked corpora \cend{50}. EigenTrust aggregates local trust scores in peer-to-peer networks with normalization to limit malicious inflation \cend{16}. TrustRank propagates trust from a seed set of known-good nodes to combat web spam \cend{17}. These systems share a theme: reputation is relational and should be hard to forge without costly endorsements from already-reputable nodes.

On-chain, Sybil attacks remain central \cend{11}. Creating wallets is cheap; creating wallets that receive allowances or transfers from diverse counterparties is costlier. Sybil-adjusted stability proxies in this research (inbound counterparty ratio, tenure, active months) partially capture diversity of interaction but do not eliminate adversarial manipulation. Chapter~\ref{ch:discussion} sketches a cost model for allowance-graph Sybils.

Machine-learning credit models on DeFi features \cend{14} can outperform simple graph scores on labeled defaults when labels exist, but they trade interpretability for accuracy. EndorseRank and AWP prioritize transparent propagation rules suitable for audit \cend{2}.

\dissertationSubheading[sec:allowance-endorsement]{ERC-20 allowances as endorsement}

ERC-20 \texttt{approve} and EIP-2612 \texttt{permit} record explicit spending authorization \cend{18,19}. EndorseRank treats the latest allowance per owner--spender pair as edge strength on a directed endorsement graph and runs PageRank on that snapshot, eliminating historical time-decay because newer approvals override older ones.

Formally, for each token and owner--spender pair, let $a_{o\to s}$ be the decimal-adjusted latest allowance. Edges are aggregated across tokens by summing allowances into a single weight $w_{o\to s}$. The resulting graph $G_{\text{approve}}=(V,E_{\text{approve}},w)$ is typically much sparser than the transfer graph over the same wallets and window. PageRank on $G_{\text{approve}}$ yields EndorseRank scores $\mathbf{s}_{\mathrm{ER}}$.

The endorsement interpretation is intentional but imperfect. Users approve routers for convenience, not always as a social judgment of the router's ``trustworthiness.'' Contract addresses and EOAs coexist. Nevertheless, allowances are a clearer authorization signal than passive receipt of transfers, which may reflect market making, wash trading \cend{29}, airdrop farming \cend{32}, or routing without any endorsement intent.

\dissertationSubheading[sub:gf-pr]{What this research does not claim}

This research does not claim that EndorseRank predicts trading profit, loss avoidance or liquidation. The research proposal planned to validate both methods against GMX V2 trading-success proxies and to bound that validation with an outcome-native reference (a PageRank on a star graph built from realized profit and loss). That reference is circular by construction: it is computed from the same PnL stream that defines the trading-success proxies, so its near-perfect agreement with the realized-gain proxy is the correlation of a variable with a monotone function of itself, not evidence about reputation. For the same reason, weak agreement between a social method and those proxies cannot be read as failure. The reference, the associated supplementary research question and claim, and the trading-success, inverse-risk and liquidation families are therefore withdrawn from the main narrative; their matrices are archived in Appendix~\ref{ch:appendix-archive} for completeness only. The remaining external check is a temporal holdout of future new owner--spender approvals (Chapters~\ref{ch:methodology} and~\ref{ch:results}), which is out of window but of the same construct as the score.

\dissertationSubheading[sec:rank-correlation-theory]{Rank correlation theory}

Reputation scores and validation proxies are continuous variables without a universal classification threshold. Following the AWP validation design\cend{3} and classical psychometrics \cend{22}, this research evaluates monotonic and pairwise rank agreement, and it quantifies the sampling uncertainty of every reported coefficient with bootstrap confidence intervals \cend{34}.

Spearman's rank correlation \cend{20} is Pearson correlation on ranks. For paired ranks $(R_i,S_i)$ of $n$ items,
\begin{equation}
\rho = 1 - \frac{6\sum_{i=1}^n (R_i-S_i)^2}{n(n^2-1)}
\label{eq:spearman}
\end{equation}
in the absence of ties (with tie-aware variants used in software). Spearman's $\rho$ captures monotonic association.

Kendall's $\tau$ \cend{21} counts concordant and discordant pairs. For all pairs $(i,j)$ with $i<j$,
\begin{equation}
\tau = \frac{2(C - D)}{n(n-1)},
\label{eq:kendall}
\end{equation}
where $C$ ($D$) is the number of concordant (discordant) pairs, with tie corrections in the $\tau$-b variant. Kendall's $\tau$ is often more interpretable for pairwise ordering agreement and is the primary family-level summary in Chapter~\ref{ch:results}. Values are compared relatively across methods on the same cohort rather than against an uncited absolute threshold.

\dissertationSubheading[sec:theoretical-framework]{Theoretical framework}

This study uses hyperlink--citation propagation, construct (domain) alignment, and computational complexity as the bases for comparing EndorseRank and AWP across reputation structure (H1), social-method alignment (H2), and computational efficiency (H3).

\begin{itemize}
\item Hyperlink--citation semantics: The original PageRank algorithm \cend{15} treats inbound hyperlinks as citations. Do and Do\cend{4} apply this logic to Ethereum transfers; AWP\cend{3} uses time-decayed transfer edge weights. Figure~\ref{fig:logistic-decay-methods} (Chapter~\ref{ch:methodology}) plots the logistic decay
\begin{equation}
\sigma(\Delta t)=\frac{1}{1+\exp\bigl(k(\Delta t-t_0)\bigr)}
\label{eq:logistic}
\end{equation}
with $k=0.01$ and $t_0=180$ days as used for AWP in this study. EndorseRank extends citation-based propagation to ERC-20 approve/permit edges because allowances are explicit, on-chain delegations of spending authority \cend{18,19}. Unlike transfers, which may reflect trading or routing without trust intent, an allowance records explicit authorization that remains in force until revoked or replaced. Graph-based risk models \cend{12} further show that network structure carries information beyond raw activity counts, which supports a directed endorsement graph in addition to transfer volume.

\item Domain alignment framework: In measurement theory, construct alignment assesses how well an indicator tracks an intended construct \cend{22}; Messick\cend{33} extends this to a unified judgment about what a score means and what its use implies. For rank-based reputation scoring, Spearman's $\rho$ and Kendall's $\tau$ evaluate monotonic and pairwise ordering agreement without assuming a classification threshold, as in the AWP validation protocol\cend{3}. Contemporaneous transfer and allowance statistics check graph coherence; a temporal holdout of future new approvals is the out-of-window check.

\item Computational complexity: PageRank's per-iteration cost is $O(|E|)$ \cend{15}. Allowance graphs are sparser than transfer graphs (fewer edges $|E|$), yielding lower per-iteration cost and reduced memory, and often fewer iterations in practice---the basis for H3. Because the solver is identical for both methods, any runtime difference is an expected consequence of edge sparsity, not an algorithmic contribution.
\end{itemize}

Figure~\ref{fig:conceptual-model} presents the conceptual model mapping the primary EndorseRank--AWP comparison to H1--H3 and RQ1--RQ4.

\begin{figure}[htbp]
\centering
% !TEX root = ../main.tex
\begin{tikzpicture}[
  font=\small,
  box/.style={draw, rounded corners, align=center, minimum width=10.5cm, minimum height=1.0cm},
  smallbox/.style={draw, rounded corners, align=center, minimum width=3.4cm, minimum height=1.05cm},
  refbox/.style={draw, dashed, rounded corners, align=center, minimum width=3.4cm, minimum height=0.95cm},
  arrow/.style={-{Stealth[length=2.5mm]}, thick}
]
\node[box] (center) at (0,3.2) {Primary comparison: EndorseRank vs AWP};

\node[smallbox] (rep) at (-5,1.2) {Reputation\\Structure};
\node[smallbox] (risk) at (0,1.2) {Construct\\Checks};
\node[smallbox] (eff) at (5,1.2) {Computational\\Efficiency};

\node[smallbox] (h1) at (-5,-1.0) {H1 / RQ1\\EndorseRank--AWP\\divergence};
\node[smallbox] (h2) at (0,-1.0) {H2 / RQ2\\Allowance vs\\transfer};
\node[smallbox] (h3) at (5,-1.0) {H3 / RQ3\\EndorseRank\\efficiency};

\node[refbox] (hold) at (0,-3.0) {RQ4: temporal holdout\\future new approvals};

\draw[arrow] (center.south -| rep.north) -- (rep.north);
\draw[arrow] (center.south) -- (risk.north);
\draw[arrow] (center.south -| eff.north) -- (eff.north);

\draw[arrow] (rep.south) -- (h1.north);
\draw[arrow] (risk.south) -- (h2.north);
\draw[arrow] (eff.south) -- (h3.north);

\draw[arrow, dashed] (h2.south) -- (hold.north);
\end{tikzpicture}

\caption{Conceptual model: primary EndorseRank vs.\ AWP comparison mapped to reputation structure (H1/RQ1), construct checks (H2/RQ2), computational efficiency (H3/RQ3), and the temporal holdout of future new approvals (RQ4).}
\label{fig:conceptual-model}
\end{figure}

\begin{itemize}
\item Holdout, not outcome ranking: H1 and H3 concern EndorseRank and AWP only. H2/RQ2 compares contemporaneous allowance and transfer checks. RQ4 is a time-split test of future new owner--spender approvals on the spender cohort.
\end{itemize}

\dissertationSubheading[sub:terminology]{Reader's guide to terminology}

The definitions in Chapter~\ref{ch:introduction} establish formal constructs; the following guide clarifies how to interpret results and avoid common misreadings:

\begin{itemize}
\item Primary methods: EndorseRank (allowance graph) and AWP (transfer graph). Both are social PageRank methods on wallet-to-wallet edges; intended as deployable reputation signals.
\item Primary graph construct: The graph a method is \emph{designed} to rank (e.g., $G_{\text{approve}}$ for EndorseRank). Distinct from checks applied after scoring.
\item Contemporaneous check: Local transfer or allowance statistics computed in the same window as the score \cend{22}.
\item Temporal holdout: Scores frozen at t1; labels from t2 new owner--spender approvals on inbound spenders.
\item Matched wallet cohort: The primary same-window sample of $n=5{,}521$ wallets with non-zero endorsement- and transfer-graph connectivity.
\item Spender holdout cohort: $n=1{,}335$ t1 inbound allowance recipients used for RQ4.
\item Expanded wallet pool: A larger address set ($N=99{,}080$) used for the node-count sensitivity of runtime beyond the matched cohort.
\end{itemize}

\dissertationSubheading[sec:domain-alignment-metrics]{Domain alignment metrics}

Following the AWP validation protocol\cend{3} and Cronbach and Meehl\cend{22}, construct alignment for rank-based scores uses Spearman's $\rho$ and Kendall's $\tau$ between reputation scores and check rankings. This research reports contemporaneous transfer and allowance checks and a temporal holdout of future new approvals; see Chapter~\ref{ch:methodology} for operational definitions.

The literature reviewed here supports three design choices: (i) PageRank as an interpretable propagation backbone; (ii) allowances as endorsement edges distinct from transfers; and (iii) rank correlation plus a time split as the check design. These choices motivate the remaining sections of this chapter and, later, the operational pipeline in Chapter~\ref{ch:methodology}.

\dissertationSubheading[sec:pagerank-worked-example]{Worked example: PageRank on a four-node endorsement graph}

To make Equation~\eqref{eq:pagerank-node} concrete, consider four wallets $\{A,B,C,D\}$ with latest-allowance edges
\[
A\to B\ (2),\quad A\to C\ (1),\quad B\to D\ (3),\quad C\to D\ (1),
\]
and damping $d=0.85$. Out-strengths are $d_A^{\mathrm{out}}=3$, $d_B^{\mathrm{out}}=3$, $d_C^{\mathrm{out}}=1$, $d_D^{\mathrm{out}}=0$ (dangling). After one power-iteration step from the uniform vector $\boldsymbol{\pi}^{(0)}=\mathbf{1}/4$, mass flows preferentially toward $D$ because both $B$ and $C$ endorse $D$, while $A$ splits mass between $B$ and $C$. Subsequent iterations mix in teleportation mass $(1-d)/4=0.0375$ at every node. The stationary ranking places $D$ highest, then $B$ and $C$, then $A$---illustrating that \emph{receiving} endorsements from nodes that themselves receive endorsements raises score, exactly the citation logic of PageRank \cend{15,23}.

If the same wallets are connected by many small transfers instead of few allowances, the transition matrix changes and the ranking can reorder even though the node set is identical. That is the microcosm of H1: EndorseRank and AWP need not agree.

\dissertationSubheading[sec:defi-risk-taxonomy]{DeFi risk taxonomy and where reputation fits}

DeFi risk is multi-layered; Werner et al.\cend{8} organise it by protocol layer, and Zhou et al.\cend{27} systematise real-world incidents by attack surface. Smart-contract risk concerns bugs and upgrade keys \cend{51}. Oracle risk concerns manipulated or stale prices \cend{39}. Market risk concerns volatility and liquidation cascades \cend{52}, including stablecoin and funding-market stress \cend{53}; the March 2020 crash illustrated how these layers interact \cend{9}. Governance risk concerns token-weighted capture \cend{54}. \emph{Counterparty} or \emph{behavioral} risk---the focus of this research---concerns whether a wallet's historical interactions suggest reliable management of leveraged exposure.

On-chain reputation scores address only the behavioral layer, and only partially. They do not replace collateral, audits, or oracle design. Their role is to summarize relational structure (who authorizes whom; who pays whom) into an ordinal signal that can be audited and recomputed \cend{2}. EndorseRank targets the authorization facet of behavioral risk; AWP targets the payment-flow facet. Neither claims to predict exploits or governance attacks.

Artificial activity further complicates behavioral measurement. Wash trading inflates volume on decentralized exchanges \cend{28} and, at scale, on unregulated centralized venues \cend{29}; front-running and other forms of extractable value generate transfers that reflect ordering games rather than relationships \cend{55}, and Qin et al.\cend{30} quantify how much of that value is realised on Ethereum. Scam tokens on Uniswap \cend{31} and cross-venue arbitrage flows \cend{56} likewise inflate transfer activity without durable counterparty relationships. Transfer graphs are especially exposed to artificial volume; allowance graphs are exposed to approval farming. Time-split checks and Sybil-oriented diagnostics are partial mitigations, not complete defenses.

\dissertationSubheading[sec:related-measurement]{Related measurement traditions}

Beyond blockchain, reputation systems have been studied in e-commerce \cend{57}, peer-to-peer networks \cend{16}, and web spam \cend{17}. Common lessons include: (i) reputation is relational; (ii) cheap identities enable Sybil attacks \cend{11}; (iii) seed sets and normalization limit inflation; and (iv) evaluation requires external criteria, not only internal consistency \cend{22}. EndorseRank inherits (i)--(iii) via PageRank damping and endorsement costs, and addresses (iv) via contemporaneous construct checks measured with Spearman's $\rho$ \cend{20} and Kendall's $\tau$ \cend{21}, and via a temporal holdout.

Personalized and topic-sensitive PageRank variants \cend{58,59} could, in future work, bias teleportation toward protocol-approved contracts or known-good wallets, analogous to TrustRank seeds. This research uses uniform teleportation for comparability with AWP\cend{3}.

\dissertationSubheading[sec:account-model]{Accounts, gas, and why logs are the measurement surface}

Ethereum-style accounts are either externally owned accounts (EOAs), controlled by private keys, or contract accounts, controlled by code \cend{60}. Every state-changing interaction consumes gas and emits logs when contracts are written to do so. Reputation research that aims for auditability therefore privileges \emph{logs} over off-chain indexes: logs are part of the canonical chain history and can be re-derived by any full node or public warehouse \cend{61}.

Gas costs shape the graphs we observe. On L2s, lower fees increase transfer frequency and approval churn, densifying $G_{\text{transfer}}$ more than $G_{\text{approve}}$ because payments are routine while approvals are intermittent setup actions. This economic asymmetry is one reason EndorseRank graphs are expected to remain sparse relative to transfer graphs (examined in Chapter~\ref{ch:results}).

\dissertationSubheading[sec:approval-security]{Approvals as security-relevant events}

Unlimited approvals (\texttt{type(uint256).max}) are a well-known user-experience hazard: a compromised spender can drain tokens until the owner revokes. Security analyses of honeypot contracts \cend{62} and of phishing accounts on Ethereum \cend{24} treat allowances as capabilities, not as neutral metadata. EndorseRank interprets inbound allowances as endorsement of a spender's capability, which is precisely why the construct is semantically distinct from inbound transfers. A wallet may receive many transfers without ever being trusted to spend on others' behalf.

From a measurement standpoint, unlimited approvals create heavy-tailed edge weights. Decimal adjustment and optional winsorization can moderate extremes; this research sums decimal-adjusted latest allowances and relies on PageRank's normalization by out-strength so that a single huge approval does not automatically dominate unless the owner has few outgoing endorsements.

\dissertationSubheading[sec:awp-in-depth]{AWP in depth: what the transfer-graph baseline measures}

Do and Do\cend{4} apply PageRank and HodgeRank to Ethereum transfers as a social-credit signal. Adaptive Weighted PageRank (AWP), introduced by Do, Do, and Nguyen\cend{3}, uses transfer edges whose weights combine value and recency through a logistic decay, and validates the ranking against inbound degree and inbound value---proxies for ``how much economic attention a wallet receives.'' Both the decay formulation (Equation~\eqref{eq:logistic}) and the in-degree/in-value proxies used in Chapters~\ref{ch:methodology} and~\ref{ch:results} come from the AWP paper. The framework is attractive because it is fully on-chain, rank-based, and free of off-chain identity.

EndorseRank preserves the PageRank-and-rank-correlation template but replaces the edge semantics. The empirical question is whether allowance edges yield a different, useful, and cheaper signal. Chapter~\ref{ch:results} tests whether allowance edges produce a construct-distinct, allowance-aligned, and computationally cheaper score, and whether t1 scores predict t2 new approvals.

\dissertationSubheading[sec:method-comparison-table]{Comparison of graph reputation methods}

Table~\ref{tab:method-comparison-lit} situates EndorseRank among related methods discussed in this chapter by edge semantics, temporal model, and validation design.

\begin{table}[htbp]
\centering
\caption{Comparison of selected graph reputation methods (conceptual).}
\label{tab:method-comparison-lit}
\fitwidth{%
\begin{tabular}{p{0.18\linewidth}p{0.24\linewidth}p{0.24\linewidth}p{0.28\linewidth}}
\toprule
Method & Edge semantics & Temporal model & Typical validation \\
\midrule
PageRank (web) \cend{15} & Hyperlinks & None / static snapshot & Retrieval quality \\
EigenTrust \cend{16} & Local trust ratings & Iterative aggregation & P2P file authenticity \\
TrustRank \cend{17} & Links from seed goods & Static / seeded & Spam demotion \\
PageRank/HodgeRank on Ethereum \cend{4} & Token transfers & Static window & Not reproduced in this thesis \\
AWP \cend{3} & Token transfers & Logistic decay & In-degree / in-value \\
RiskProp \cend{12} & Risk-bearing links & Propagation of risk & Account risk labels \\
EndorseRank (this study) & Latest allowances & Snapshot (override) & Allowance checks; t2 new approvals \\
\bottomrule
\end{tabular}
}
\end{table}

Table~\ref{tab:method-comparison-lit} shows that EndorseRank's novelty is the choice of latest-allowance endorsement edges as the primary social construct for DeFi wallets, not a new ranking algorithm.

\dissertationSubheading[sec:literature-summary-table]{Summary of the reviewed literature}

Table~\ref{tab:literature-summary} consolidates the sources that shape the design of this study. For each source it records the methodology, the strengths that this thesis borrows, the limitations that it tries to address, and the relation to the present work. The table is the evidence base for the research gap stated in Section~\ref{sec:research-gap}.

{\footnotesize
\setlength{\tabcolsep}{4pt}
\begin{longtable}{p{0.15\linewidth}p{0.20\linewidth}p{0.19\linewidth}p{0.19\linewidth}p{0.18\linewidth}}
\caption{Summary of the reviewed literature: methodology, strengths, limitations, and relation to this study.}
\label{tab:literature-summary}\\
\toprule
Source (year) & Methodology & Strengths & Limitations & Relation to this study \\
\midrule
\endfirsthead
\multicolumn{5}{l}{\small Table~\thetable\ (continued)}\\
\toprule
Source (year) & Methodology & Strengths & Limitations & Relation to this study \\
\midrule
\endhead
\midrule
\multicolumn{5}{r}{\small continued on next page}\\
\endfoot
\bottomrule
\endlastfoot
Page et al.\ (1998) \cend{15}; Langville and Meyer (2006) \cend{23} & Eigenvector ranking of a hyperlink graph; power iteration & Interpretable propagation; $O(|E|)$ per iteration; unique Perron vector under damping & Edge semantics fixed to hyperlinks; no notion of time & Solver and complexity argument reused unchanged for both compared methods \\
Kamvar et al.\ (2003) \cend{16}; Gy\"{o}ngyi et al.\ (2004) \cend{17} & Trust aggregation in P2P networks; seeded propagation against spam & Normalisation and seed sets limit inflation; explicit adversary model & Requires explicit ratings or curated seeds; not on-chain & Motivates the Sybil-cost discussion and the uniform-teleportation baseline \\
Do and Do (2023) \cend{4} & PageRank and HodgeRank on Ethereum transaction graphs as social credit & Framing of transfer connectivity as credit-relevant; fully public data & Transfers conflate payment and trust; no allowance semantics & Prior work on the transfer graph; distinguished from AWP in this thesis \\
Do, Do, and Nguyen (2023), AWP \cend{3} & Value- and recency-weighted transfer PageRank with logistic decay; rank correlation against in-degree/in-value & Rank-based validation without default labels; recency handled explicitly & Proxies are computed from the same transfer data that builds the graph (mechanical alignment); dense graph & Baseline method, decay formula, and proxy protocol reproduced exactly on the same cohort \\
Lin et al.\ (2025), RiskProp \cend{12} & De-anonymisation score plus network propagation of account risk & Shows structure carries information beyond activity counts; labelled evaluation & Needs risk labels; targets fraud, not trust or authorisation & Motivates refined edge semantics; RiskProp-style propagation kept as an archived extended baseline \\
Wu et al.\ (2022) \cend{24} & Network embedding of Ethereum transaction graphs for phishing detection & Graph features outperform per-account features; large labelled set & Supervised; labels are phishing, not creditworthiness & Supports the premise that graph position is informative; not a competitor method \\
Hassija et al.\ (2020) \cend{13}; Kotb (2023) \cend{14} & Prospect-theory and machine-learning credit scoring on blockchain data & Directly target lending decisions & Require labels or off-chain features; limited interpretability & Contrast class: EndorseRank stays within auditable graph propagation \\
Bellini et al.\ (2020) \cend{25}; Almasoud et al.\ (2020) \cend{26} & Surveys of blockchain-based trust and reputation systems & Map the design space; identify explicit-rating and volume-based signals as dominant & Do not evaluate authorisation edges; no L2 measurement & Establish that latest-allowance edges are absent from the surveyed designs \\
Victor (2020) \cend{40}; Victor and Weintraud (2021) \cend{28}; Cong et al.\ (2023) \cend{29} & Address clustering heuristics; wash-trading detection on DEXs and exchanges & Quantify how much on-chain activity is artificial or multi-address & Do not propose a reputation score & Explain why transfer-based signals are exposed to inflation; motivate the allowance construct \\
Cronbach and Meehl (1955) \cend{22}; Messick (1995) \cend{33} & Construct validity in psychometrics & Separate what a score is designed to measure from external criteria & Not blockchain-specific & Frame the contemporaneous checks as construct checks and the temporal holdout as external evidence \\
Efron (1979) \cend{34}; DiCiccio and Efron (1996) \cend{35} & Bootstrap resampling and bootstrap confidence intervals & Distribution-free uncertainty for rank statistics; paired contrasts & Assumes exchangeable units & Provides the 95\% intervals and $\Delta\tau$ contrasts reported in Chapter~\ref{ch:results} \\
\end{longtable}
}

\dissertationSubheading[sec:research-gap]{Research gap}

Three gaps follow from Table~\ref{tab:literature-summary}. First, every on-chain reputation method reviewed---PageRank/HodgeRank \cend{4}, AWP \cend{3}, RiskProp \cend{12}, and the systems catalogued in the surveys \cend{25}---builds its graph from transfers or explicit ratings; none uses the ERC-20 allowance, the only on-chain event that records a deliberate delegation of spending authority. Whether a latest-allowance endorsement graph produces a different and useful ordering is therefore untested (addressed by O1 and RQ1). Second, the validation protocol of the transfer-graph baseline compares scores with in-degree and in-value computed from the same transfer data that builds the graph, so agreement is partly mechanical; no study reports the same protocol on an allowance construct alongside its transfer counterpart with confidence intervals and paired difference tests on one cohort (O2 and RQ2). Third, reported runtimes for blockchain PageRank are rarely accompanied by edge counts, memory, and iteration counts under a fixed protocol, which makes efficiency claims hard to attribute to sparsity versus implementation (O3 and RQ3). A fourth, methodological gap concerns time: same-window graph statistics are routinely presented as validation even though they cannot show that a score anticipates future behaviour; a temporal holdout on future new approvals is the out-of-window check adopted here (O4 and RQ4).

\dissertationSubheading[sec:open-problems]{Open problems motivating this research}

Several open problems in on-chain reputation remain unresolved in the literature and motivate the present study:

\begin{enumerate}
\item Edge semantics: Which on-chain relations best encode trust or authorization under pseudonymity \cend{2}?
\item Validation without defaults: How can methods be compared when loan-default labels are scarce \cend{4}?
\item Computational cost: Can reputation scores refresh at interactive latencies as graphs grow \cend{23}?
\item Time-split checks: How can evaluation avoid mistaking same-window graph statistics for out-of-window evidence \cend{22}?
\item Adversarial robustness: What is the cost of forging reputation under cheap identities \cend{11}?
\end{enumerate}

EndorseRank addresses (1)--(4) with allowance edges, contemporaneous checks, runtime benchmarks, and a temporal holdout, and analyzes (5) at a model level in Chapter~\ref{ch:discussion}. Full empirical adversarial evaluation remains future work.

Chapter~\ref{ch:methodology} turns the gap statement into an operational design: it fixes the data window and cohort, specifies how both graphs are constructed and solved with one shared PageRank implementation, defines the proxy families and the temporal holdout, and sets the timing and bootstrap protocols whose outputs Chapter~\ref{ch:results} reports.

